\documentclass[11pt]{article}
\usepackage[margin=1in]{geometry}
\usepackage{amsmath,amssymb,amsthm}
\usepackage{hyperref}

\theoremstyle{plain}
\newtheorem{theorem}{Theorem}[section]
\newtheorem{proposition}[theorem]{Proposition}
\newtheorem{lemma}[theorem]{Lemma}
\newtheorem{corollary}[theorem]{Corollary}
\theoremstyle{definition}
\newtheorem{definition}[theorem]{Definition}
\theoremstyle{remark}
\newtheorem{remark}[theorem]{Remark}

\title{Notes on the Riemann Hypothesis}
\author{KernelCert Workspace Notes}
\date{August 10, 2026}

\begin{document}
\maketitle

\section*{Important Note}
The Riemann Hypothesis (RH) is an open problem in mathematics as of the date above.
No accepted proof is known.
Accordingly, this document does \emph{not} present a proof; it records a precise statement,
standard equivalent viewpoints, and where known approaches face major obstacles.

\section{Statement of the Problem}
\begin{definition}[Riemann zeta function]
For $\Re(s)>1$, define
\[
\zeta(s) = \sum_{n=1}^{\infty} \frac{1}{n^s}.
\]
It has a meromorphic continuation to $\mathbb{C}$ with a unique simple pole at $s=1$.
\end{definition}

\begin{lemma}[Location of nontrivial zeros]
All nontrivial zeros of $\zeta(s)$ lie in the critical strip
\[
0<\Re(s)<1.
\]
\end{lemma}

\begin{proof}[Proof sketch]
Use the Euler product to rule out zeros on $\Re(s)>1$, then apply the functional equation and
symmetry to transport this exclusion to $\Re(s)<0$ except for trivial zeros at negative even integers.
\end{proof}

\begin{theorem}[Riemann Hypothesis (open)]
Every nontrivial zero $\rho$ of $\zeta(s)$ satisfies
\[
\Re(\rho)=\tfrac12.
\]
\end{theorem}

\begin{remark}
The theorem above is currently unproved; in modern usage it is a conjecture.
\end{remark}

\section{Equivalent and Near-Equivalent Perspectives}
The next statements are classical equivalences to RH and are included as formal targets.

\begin{proposition}[Prime counting error formulation]
Let $\pi(x)$ be the prime counting function and $\operatorname{Li}(x)$ the logarithmic integral.
Then RH is equivalent to
\[
\pi(x)=\operatorname{Li}(x)+O\!\left(\sqrt{x}\log x\right).
\]
\end{proposition}

\begin{proof}[Proof sketch]
Apply explicit formulas linking primes and zeros of $\zeta$, then bound contributions from zeros.
The direction RH $\Rightarrow$ error term uses zero locations; the converse uses contour estimates to
exclude zeros off the critical line.
\end{proof}

\begin{proposition}[Chebyshev function formulation]
For
\[
\psi(x)=\sum_{p^k\le x}\log p,
\]
RH is equivalent to
\[
\psi(x)=x+O\!\left(\sqrt{x}(\log x)^2\right).
\]
\end{proposition}

\begin{proof}[Proof sketch]
Use von Mangoldt's explicit formula. The oscillatory term is governed by nontrivial zeros; the stated
error term is exactly the threshold corresponding to real part $1/2$.
\end{proof}

\begin{proposition}[Li criterion]
Let $\xi(s)$ be Riemann's xi-function and define
\[
\lambda_n = \frac{1}{(n-1)!}\left.\frac{d^n}{ds^n}\left(s^{n-1}\log\xi(s)\right)\right|_{s=1}
\quad (n\ge 1).
\]
Then RH is equivalent to $\lambda_n\ge 0$ for all $n\ge 1$.
\end{proposition}

\begin{proof}[Proof sketch]
Express $\log\xi$ through its zero set and identify $\lambda_n$ as weighted sums over zeros.
Nonnegativity for all $n$ matches the critical-line condition via Li's theorem.
\end{proof}

\section{First Crack: Potential Solution Paths (Not a Proof)}
This section recasts potential approaches as formal research propositions.
Each item is intentionally marked as incomplete.

\begin{proposition}[Program A: Hilbert-Polya spectral route]
Assume there exists a densely defined self-adjoint operator $T$ with a trace formula matching the
zeta explicit formula and whose spectrum corresponds to the imaginary parts of nontrivial zeros.
Then RH follows.
\end{proposition}

\begin{proof}[Proof sketch]
Self-adjointness forces real spectral parameters. The spectral identification then places all
nontrivial zeros on the critical line after normalization.
\end{proof}

\begin{remark}[Status]
No construction of such an operator is known with all required analytic and spectral properties.
\end{remark}

\begin{proposition}[Program B: positivity from explicit-formula kernels]
Assume one can build a test-function class for Weil's explicit formula such that any zero with
$\Re(\rho)\ne 1/2$ contributes a strictly sign-definite term violating global positivity.
Then RH follows.
\end{proposition}

\begin{proof}[Proof sketch]
Under the stated kernel properties, existence of an off-line zero implies a sign contradiction in the
explicit formula, forcing all zeros to satisfy $\Re(\rho)=1/2$.
\end{proof}

\begin{remark}[Status]
Known kernel-positivity frameworks are not yet strong enough to enforce this global contradiction.
\end{remark}

\begin{proposition}[Program C: de Bruijn-Newman route]
Let $\Lambda$ denote the de Bruijn-Newman constant. If one proves $\Lambda\le 0$, then RH follows.
\end{proposition}

\begin{proof}[Proof sketch]
Classically, RH is equivalent to $\Lambda\le 0$ for the heat-flow deformation of $\xi$; thus the claim is immediate.
\end{proof}

\begin{remark}[Status]
Available lower bounds and computations indicate proximity to $0$ but do not prove $\Lambda\le 0$.
\end{remark}

\begin{proposition}[Program D: mollifier and moment route]
Assume one can prove moment estimates and zero-density bounds strong enough to exclude every off-line
zero after mollification. Then RH follows.
\end{proposition}

\begin{proof}[Proof sketch]
The hypothesis collapses off-line zero density to zero while retaining all nontrivial zeros in the strip,
hence every nontrivial zero must lie on $\Re(s)=1/2$.
\end{proof}

\begin{remark}[Status]
Current methods establish substantial partial results (including large proportions on the line),
but not complete exclusion of off-line zeros.
\end{remark}

\begin{lemma}[Minimal completion criteria]
Any complete proof attempt should provide:
\begin{enumerate}
  \item a mechanism forcing all nontrivial zeros to have real part $1/2$;
  \item a non-circular derivation avoiding RH-equivalent assumptions;
  \item full analytic error control suitable for independent verification.
\end{enumerate}
\end{lemma}

\begin{remark}
No path in this document currently satisfies all items of the lemma.
\end{remark}

\section{Candidate Breakthrough Skeleton (Conditional)}
This section records a tighter theorem-style blueprint that would look like a genuine breakthrough
\emph{if} all hypotheses were proved.

\begin{definition}[Weil quadratic form on test functions]
Let $\mathcal{T}$ denote an admissible class of even test functions used in Weil's explicit formula.
Write $Q(f)$ for the corresponding zero-side quadratic form.
\end{definition}

\begin{lemma}[Rigidity principle: conditional form]
Assume there exists a subclass $\mathcal{T}_{\mathrm{rigid}}\subseteq\mathcal{T}$ and constants
$c_0>0$, $R_0\ge 1$ such that:
\begin{enumerate}
  \item $Q(f)\ge 0$ for every $f\in\mathcal{T}_{\mathrm{rigid}}$;
  \item for any zero $\rho=\beta+i\gamma$ with $\beta\ne 1/2$ and $|\gamma|\ge R_0$, there exists
  $f_\rho\in\mathcal{T}_{\mathrm{rigid}}$ such that the off-line contribution is at most $-c_0$;
  \item there exists a finite verification bound for $|\gamma|<R_0$ excluding off-line zeros in that range.
\end{enumerate}
Then RH holds.
\end{lemma}

\begin{proof}[Proof sketch]
An off-line zero with $|\gamma|\ge R_0$ contradicts item (1) via item (2). The range
$|\gamma|<R_0$ is handled by item (3). Hence all nontrivial zeros satisfy $\Re(\rho)=1/2$.
\end{proof}

\begin{remark}[Gap A]
No currently known class $\mathcal{T}_{\mathrm{rigid}}$ is proved to satisfy both items globally.
This is the key unresolved bottleneck in the positivity route.
\end{remark}

\begin{proposition}[Density-to-rigidity transfer: conditional]
Assume there exist constants $A,B>0$, $T_0\ge 3$, and $\eta\in(0,1/2)$ such that:
\begin{enumerate}
  \item for all $T\ge T_0$ and all $\sigma\in[1/2+\eta,1)$,
\[
N(\sigma,T) \le A\,T^{2(1-\sigma)}(\log T)^B
\]
  where $N(\sigma,T)$ counts zeros with $\Re(\rho)\ge\sigma$ and $|\Im(\rho)|\le T$;
  \item there exist $C_1,C_2>0$ such that for all $f\in\mathcal{T}_{\mathrm{rigid}}$ and all $T\ge T_0$,
  the explicit-formula remainder term $\mathcal{E}(f;T)$ satisfies
\[
|\mathcal{E}(f;T)| \le C_1\,\|f\|_{\ast}(\log T)^{C_2},
\]
  for an admissible norm $\|\cdot\|_{\ast}$ on test functions.
\end{enumerate}
Then any off-line zero forces a contradiction with the stability estimate.
\end{proposition}

\begin{proof}[Proof sketch]
The density bound controls how frequently off-line zeros can appear at scale $T$.
The stability estimate amplifies one hypothetical off-line zero into a sign defect in the explicit formula,
incompatible with global positivity.
\end{proof}

\begin{remark}[Gap B]
The needed stability estimate is currently stronger than what existing mollifier and moment methods deliver.
\end{remark}

\begin{proposition}[Spectral closure: conditional]
Suppose Program A provides a self-adjoint operator model with trace identity that matches the prime-side
terms of Weil's explicit formula for all $f\in\mathcal{T}_{\mathrm{rigid}}$, and suppose the model is
defined on a dense domain $\mathcal{D}(T)$ with essential self-adjointness and spectral measure
having polynomial growth in frequency.
Then the rigidity lemma above is consistent with a spectral proof template.
\end{proposition}

\begin{proof}[Proof sketch]
The trace identity identifies the same distributional object from both spectral and arithmetic sides.
Self-adjointness enforces real spectral parameters, which aligns with the critical-line condition.
\end{proof}

\begin{remark}[Gap C]
Constructing the operator and proving its full domain/spectral properties remain open.
\end{remark}

\begin{theorem}[Conditional breakthrough theorem]
If Gap A, Gap B, and Gap C are all resolved with fully rigorous proofs, then RH follows.
\end{theorem}

\begin{proof}
Combine the rigidity lemma, density-to-rigidity transfer, and spectral closure proposition.
Each removes one failure mode in the current landscape; jointly they force all nontrivial zeros onto
$\Re(s)=1/2$.
\end{proof}

\begin{remark}
This theorem is a roadmap statement, not a completed proof.
It is the closest part of this draft to a ``breakthrough shape,'' but decisive hypotheses are still open.
\end{remark}

\subsection*{Milestones and Proof Obligations}
The following checklist records concrete obligations needed to close each named gap.

\begin{proposition}[Milestone A (for Gap A)]
Construct an explicit family $\{f_{\theta}\}_{\theta\in\Theta}\subseteq\mathcal{T}_{\mathrm{rigid}}$ and
prove uniformly that:
\begin{enumerate}
  \item $Q(f_{\theta})\ge 0$ for all $\theta\in\Theta$;
  \item for each off-line zero $\rho$ with $|\Im(\rho)|\ge R_0$, at least one $f_{\theta}$ detects it with
  negativity margin at least $c_0$.
\end{enumerate}
\end{proposition}

\begin{proposition}[Milestone B (for Gap B)]
Prove the remainder bound
\[
|\mathcal{E}(f;T)| \le C_1\,\|f\|_{\ast}(\log T)^{C_2}
\]
uniformly over $f\in\mathcal{T}_{\mathrm{rigid}}$ and $T\ge T_0$, and establish a zero-density estimate
valid uniformly on $\sigma\in[1/2+\eta,1)$.
\end{proposition}

\begin{proposition}[Milestone C (for Gap C)]
Produce an operator model $(T,\mathcal{D}(T))$ with:
\begin{enumerate}
  \item essential self-adjointness on a dense domain;
  \item a trace identity matching the arithmetic explicit formula on $\mathcal{T}_{\mathrm{rigid}}$;
  \item spectral growth bounds compatible with the stability constants in Milestone B.
\end{enumerate}
\end{proposition}

\begin{remark}
Milestones A--C are designed so that their conjunction implies the hypotheses of the
conditional breakthrough theorem.
\end{remark}

\section{Non-Completion Declaration}
\begin{theorem}[What is and is not proved here]
This document does not prove RH.
\end{theorem}

\begin{proof}
All theorem-level claims beyond classical equivalences are conditional research programs, and each is
accompanied by an explicit unresolved obstacle.
\end{proof}

\section{Build Note}
This file is written to be conservative for \texttt{pdflatex} (standard article class and AMS packages only).
To compile locally when a TeX distribution is installed, run:
\begin{verbatim}
pdflatex -interaction=nonstopmode -halt-on-error riemann_hypothesis.tex
\end{verbatim}

\section{References (Starting Points)}
\begin{itemize}
  \item H. M. Edwards, \emph{Riemann's Zeta Function}.
  \item E. C. Titchmarsh and D. R. Heath-Brown, \emph{The Theory of the Riemann Zeta-Function}.
  \item A. Ivic, \emph{The Riemann Zeta-Function}.
  \item J. B. Conrey, ``The Riemann Hypothesis,'' \emph{Notices of the AMS} 50 (2003), 341--353.
\end{itemize}

\end{document}

